\clearpage
\SourceLeaf{415}%
\begin{center}\large NEUE ALLGEMEINE UNTERSUCHUNGEN\\[4pt] ÜBER\\[4pt] DIE KRUMMEN FLÄCHEN.\\[7pt] [1825.]\end{center}
\noindent\hspace*{1.2em}Wenn gleich der Zweck dieses Werks eigentlich in der Ausführung neuer seinen Gegenstand betreffender Lehren besteht, so werden wir doch anfangs auch das schon Bekannte entwickeln, theils des Zusammenhanges und der Vollständigkeit wegen, theils weil die Art unserer Behandlung von der bisherigen verschieden ist. Wir werden sogar zuerst einiges die Curven in der Ebene betreffende aus denselben Gründen vorausschicken.

\begin{center}1.\end{center}
\noindent\hspace*{1.2em}Um die verschiedenen Richtungen von geraden Linien in Einer Ebene bequem unter einander vergleichen zu können, denke man sich in derselben Ebene um einen beliebigen Mittelpunkt einen Kreis mit dem Halbmesser $1$ beschrieben; die Lage des mit einer vorgegebenen Linie parallel gezogenen Radius dieses Kreises versinnlicht dann die Lage von jener, und die Neigung zweier geraden Linien gegen einander wird durch den Winkel der beiden repräsentirenden Radien, oder durch den Bogen zwischen ihren Endpunkten gemessen. Es versteht sich, dass, wo genaue Bestimmung nöthig ist, bei jeder
